\begin{graphicspathcontext}{{./chapters/ai/examples/imgs/auto/},{./chapters/ai/examples/imgs/},\old}

\begin{frame}[t]{{Les hallucinations des IA} : un LLM ne raisonnent pas au sens strict}
	\smaller\smaller
	\alertbox*{Comment expliquer la précaution des pompiers ici ?}
	\begin{columns}[T]
		\begin{column}{.5\linewidth}
			\begin{block}{ChatGPT -- 25/05/2026}
				Les pompiers ont placé des protections rouges sur les rails pour faire passer les tuyaux sans les écraser lorsqu’un train traverse le passage à niveau.
				Un train est extrêmement lourd : si un tuyau d’incendie reste directement sur le rail, il peut :
				\begin{itemize}
				\item être écrasé ou coupé,
				\item perdre la pression d’eau,
				\item interrompre l’alimentation des lances pendant l’intervention.
				\end{itemize}
				Les plaques rouges servent donc de ``ponts de franchissement'' :
				\begin{itemize}
				\item elles répartissent le poids du train,
				\item maintiennent un espace pour le tuyau,
				\item évitent aussi que le tuyau ne se coince dans les rails.
				\end{itemize}
			\end{block}
		\end{column}
		\begin{column}{.5\linewidth}
			\centering\includegraphics[width=.8\linewidth]{pb_chatgpt}
			\begin{block}{\insertcontinuationtext}
				Cette précaution permet aux secours de continuer à pomper de l’eau même si la voie ferrée doit rester ouverte à la circulation.
			\end{block}
		\end{column}
	\end{columns}
\end{frame}

\end{graphicspathcontext}

\endinput
